% Part: computability
% Chapter: computability-theory
% Section: def-functions-self-reference

\documentclass[../../../include/open-logic-section]{subfiles}

\begin{document}

\olfileid{cmp}{thy}{slf}
\olsection{تعريف الدوال بالإحالة الذاتية}

من المفيد بوجه عام أن نستطيع تعريف الدوال بدلالة نفسها. فمثلًا، إذا
أعطيت دوالًا قابلة للحساب $k$ و$l$ و~$m$، أفادت لمة النقطة الثابتة
بوجود دالة جزئية قابلة للحساب~$f$ تحقق المعادلة الآتية لكل~$y$:
\[
f(y) \simeq
\begin{cases}
k(y) & \text{إذا كان $l(y) = 0$} \\
f(m(y)) & \text{خلاف ذلك.}
\end{cases}
\]
وبصياغة أدق، تنتج $f$ بأن نضع
\[
g(x,y) \simeq
\begin{cases}
k(y) & \text{إذا كان $l(y) = 0$} \\
\cfind{x}(m(y)) & \text{خلاف ذلك}
\end{cases}
\]
ثم نستعمل لمة النقطة الثابتة لإيجاد فهرس~$e$ بحيث
$\cfind{e}(y)\simeq g(e,y)$.

وكمثال عيني يمكن تعريف دالة «القاسم المشترك الأكبر»
$\fn{gcd}(u,v)$ بالمعادلة
\[
\fn{gcd}(u,v) \simeq
\begin{cases}
v & \text{إذا كان $u = 0$} \\
\fn{gcd}(\fn{mod}(v, u), u) & \text{خلاف ذلك}
\end{cases}
\]
حيث تدل $\fn{mod}(v,u)$ على باقي قسمة $v$ على~$u$. ويبين تطبيق لمة
النقطة الثابتة أن $\fn{gcd}$ جزئية قابلة للحساب. (ويمكن في الواقع
صب هذا في القالب السابق بجعل $y$ ترمز الزوج $\tuple{u,v}$.) ثم يبين
استقراء لاحق على~$u$ أن $\fn{gcd}$ كلية في الواقع.

وبالطبع يمكن إنشاء تعريفات ذاتية الإحالة أعقد كثيرًا من الأمثلة التي
ناقشناها للتو. وتتيح معظم لغات البرمجة، بطريقة أو بأخرى، تعريف الدوال
بدلالة نفسها. ولاحظ أن هذا أقل إثارة قليلًا من القدرة على تعريف دالة
بدلالة \emph{فهرس} لخوارزمية تحسب الدالة، وهو ما تتيحه مبرهنة النقطة
الثابتة في كامل عمومها.

\end{document}
